\documentclass{article}



\usepackage{graphicx}
\usepackage{setspace}
\usepackage{amsfonts}
\usepackage{amsmath}

\usepackage{fancyhdr}
\usepackage{enumitem}
\usepackage{indentfirst}
\usepackage{glossaries}
\usepackage{lipsum}
\usepackage{xepersian}
\settextfont[Scale=1.2]{B Nazanin}
\setlatintextfont[Scale=1]{Times New Roman}
\defpersianfont\Titr[Scale=3]{B Titr}
\usepackage{xepersian}

\begin{document}


محاسبات کوانتومی بر پایه سیستم های مکانیک  کوانتومی صورت می پذیرد. مکانیک کوانتومی چارچوبی ریاضی یا مجموعه ای از قوانین برای ساخت تئوری های فیزیکی می باشد.
واحد پایه ای در یک کامپیوتر کلاسیک بیت می باشد که می تواند دو مقدار 0 یا 1 را داشته باشد. در یک روش مشابه برای محاسبات کوانتومی واحد پایه ای اطلاعات کیوبیت می باشد که مخفف عبارت بیت کوانتومی بوده و پایه اطلاعات در یک سیستم کوانتومی دو سطحی می باشد. بر خلاف حالت کلاسیک که یک بیت می تواند دارای مقدار صفر یا یک باشد، یک کیوبیت می تواند در یک ترکیب یا ابر برهم نهی از این دو حالت باشد. نحوه ی نمایش یک کیوبیت بصورت رابطه زیر نشان داده می شود:




\end{document}